% Gerado por app/nlp/training/train_toxicity_llm_baselines.py
% (ao fim da rodada, ou com --so-apendice). NAO EDITAR A MAO.
% Rodada: video-v2 | resultados: llm_baselines_toxicity_20261120_161717.json
% Linhas longas do prompt quebradas em 88 caracteres; continuacoes com 2 espacos.
% O verbatim vai num figure*: no template de duas colunas do JBCS a coluna tem
% 240pt (~56 caracteres em \footnotesize), e a largura total, 498pt (~117).
\begin{figure*}[!t]
\begingroup\footnotesize
\begin{verbatim}
[system]
Você é um sistema de moderação de conteúdo em português brasileiro. Sua tarefa é
  classificar o grau de toxicidade de uma frase extraída da transcrição de um vídeo, em
  três níveis:
0 = Nenhuma: texto apropriado, isento de linguagem agressiva.
1 = Leve: ofensividade indireta — críticas dirigidas a objetos, situações ou entidades
  abstratas, sem alvo pessoal (ex.: "Que porra de ônibus").
2 = Severa: ofensividade direta — insultos, ameaças ou discurso de ódio dirigidos
  explicitamente a uma pessoa ou grupo (ex.: "Este idiota do Robson fez isso").
Responda APENAS com o número 0, 1 ou 2, sem nenhuma outra palavra.

[user]
Frase: lanchinhos comprados e agora livros

[assistant]
0

[user]
Frase: Cala essa boca, seu criolo safado! Você não tem o menor respeito por nada!

[assistant]
2

[user]
Frase: seu puto fudido

[assistant]
2

[user]
Frase: Vai a merda

[assistant]
1

[user]
Frase: Por que essa merda?

[assistant]
1

[user]
Frase: eu perguntei pro Léo ah o Léo disse que gosta de brincar com as fantasias ótimo
  vamos Léo

[assistant]
0

[user]
Frase: <sentence to classify>
\end{verbatim}
\endgroup
\caption{Messages sent to the generative baseline for each test sentence, in order; the
role of each message is given in brackets. Lines longer than 88 characters are wrapped
here for typesetting, with continuation lines indented by two spaces; blank lines separate
messages, and every other line break in the listing is part of the prompt. The final user
message carries the sentence being classified.}
\label{fig:prompt-llm}
\end{figure*}

\begin{itemize}
\item Model: \texttt{llama3.1:8b}, served locally by Ollama through the \texttt{/api/chat}
  endpoint, one request per sentence, without streaming.
\item Decoding: \texttt{temperature} = 0, \texttt{num\_predict} = 5; all other options at
  the server defaults.
\item Few-shot examples: 2 per class, drawn from the training split with a fixed seed
  (42); the test split is never used.
\item Parsing: if the stripped reply is exactly \texttt{0}, \texttt{1} or \texttt{2}, that
  is the label; otherwise the first of these characters anywhere in the reply; otherwise
  the first keyword found, tested in the order \emph{severa}, \emph{leve}, \emph{nenhuma};
  otherwise the sentence is labelled None. No sentence is discarded or re-sent.
\item Replies outside the requested format: \texttt{llama3.1:8b}: 358 of 358 replies in
  the requested format, 0 recovered from a digit, 0 from a keyword, and 0 with no
  recognisable label (assigned None).
\end{itemize}
